\documentclass[12pt]{article}
\usepackage[T1]{fontenc}
\usepackage[utf8]{inputenc}
\usepackage{lmodern}
\usepackage{textcomp}
\usepackage{enumitem}
\usepackage{geometry}
\geometry{margin=1in}
\begin{document}
\begin{center}
Answer Key - Economics - Undergraduate Level Assessment 10 \
\small (Answers are ordered exactly as in the question paper.)
\end{center}

\section*{Multiple Choice}
\begin{enumerate}[label=\arabic*.]
\item \textbf{Answer:} A
\\ \textit{Options:} A) Capital, commodities, and currency; B) raw materials, human resources and technology; C) natural resources, artificial resources and human resources; D) Water, oil, and minerals; E) Intellectual property, workforce, and real estate
\\ \textit{Note:} The correct answer identifies capital, commodities, and currency as essential categories of economic resources or factors of production, reflecting the fundamental components required for the production of goods and services in an economy, where capital represents tools and infrastructure, commodities are raw materials, and currency facilitates transactions.
\item \textbf{Answer:} A
\\ \textit{Options:} A) tariffs protecting domestic jobs are eliminated.; B) the literacy rate improves.; C) the tax rate is lowered.; D) the inflation rate decreases.; E) there is an increase in technological advancements.
\\ \textit{Note:} Eliminating tariffs that protect domestic jobs can lead to increased competition and efficiency among firms, resulting in a reallocation of resources that may reduce potential GDP in the short term as industries adjust to a more open market.
\item \textbf{Answer:} A
\\ \textit{Options:} A) The Antitrust Improvement Act; B) The Federal Trade Commission Act; C) The Clayton Act; D) The Sherman Act; E) The Fair Trade Act
\\ \textit{Note:} The Antitrust Improvement Act is not considered a landmark antitrust act because it primarily served to amend existing legislation rather than establish foundational antitrust principles like the Sherman Act or the Clayton Act.
\item \textbf{Answer:} D
\\ \textit{Options:} A) Capitalism results in single authority; B) Capitalism ensures job security; C) Capitalism fosters cultural diversity; D) Capitalism promotes personal freedom; E) Capitalism encourages moral responsibility
\\ \textit{Note:} The correct answer is rooted in the principle that capitalism allows individuals to make their own choices regarding production, consumption, and investment, thereby fostering a sense of personal autonomy and freedom that is not as prevalent in other economic systems.
\item \textbf{Answer:} A
\\ \textit{Options:} A) stabilizes the value of currency over time.; B) has no impact on purchasing power.; C) increases the real value of debt.; D) encourages households to save more.; E) forces households to save more.
\\ \textit{Note:} Inflation, when moderate and stable, helps to prevent deflation and encourages spending and investment, which in turn stabilizes the currency's value over time by maintaining consumer confidence and economic growth.
\end{enumerate}

\section*{True / False}
\begin{enumerate}[label=\arabic*.]
\setcounter{enumi}{5}
\item \textbf{Answer:} False
\\ \textit{Options:} A) True; B) False
\\ \textit{Note:} The statement is false because when the market price is above the average total cost, firms are incentivized to enter the market to take advantage of the economic profits, leading to an increase in industry supply rather than an unchanged market.
\item \textbf{Answer:} False
\\ \textit{Options:} A) True; B) False
\\ \textit{Note:} The statement is false because a nation at the natural rate of unemployment is not necessarily experiencing significant inflation, as this rate reflects a balance between employment and inflationary pressures, typically associated with stable prices rather than high inflation.
\item \textbf{Answer:} True
\\ \textit{Options:} A) True; B) False
\\ \textit{Note:} The demand curve for labor is derived from the demand curve for the output because employers hire labor based on the value generated from the output produced, meaning as the demand for the final product increases, so does the demand for the labor needed to produce it.
\item \textbf{Answer:} False
\\ \textit{Options:} A) True; B) False
\\ \textit{Note:} The statement is false because according to Keynesian theory, a decrease in the money supply actually raises interest rates, which discourages borrowing and reduces spending, contrary to what is claimed.
\item \textbf{Answer:} False
\\ \textit{Options:} A) True; B) False
\\ \textit{Note:} The statement is false because the bank's actual reserves are the sum of required reserves and excess reserves, but it is incorrect to claim that they equal \$50,000 without specific information about the actual figures for required and excess reserves.
\end{enumerate}

\section*{Long Form Response}
\begin{enumerate}[label=\arabic*.]
\setcounter{enumi}{10}
\item \textbf{Answer:} Product differentiation plays a crucial role in monopolistic competition by allowing each firm to distinguish its products from those of competitors, thereby giving them some control over pricing. This differentiation can be based on various factors such as quality, design, branding, or customer service. As a result, firms are not price takers like in perfect competition; instead, they can set prices above marginal cost due to the perceived uniqueness of their offerings. This ability to influence prices leads to a variety of market dynamics, including increased competition for consumer loyalty and innovation, as firms strive to enhance their product features and marketing strategies. Ultimately, product differentiation not only shapes individual firm behavior but also influences overall market structure and consumer choices.
\\ \textit{Note:} correct because it accurately describes how product differentiation in monopolistic competition enables firms to exert pricing power by creating perceived uniqueness in their products, allowing them to set prices above marginal cost unlike price-taker firms in perfect competition.
\item \textbf{Answer:} A progressive income tax system contributes to economic stability by introducing a built- in stabilizer that adjusts tax burdens based on individual income levels. During economic upswings, higher earners contribute more, which can help moderate inflation and prevent overheating in the economy. Conversely, in downturns, the tax burden on lower and middle- income earners decreases, leaving them with more disposable income to sustain consumer spending, thus aiding in economic recovery. Additionally, the progressive tax structure ensures that government revenue remains robust during prosperous times while providing necessary fiscal flexibility during recessions, thereby supporting essential public services and stabilizing the economy. Overall, this system fosters a more resilient economic environment by balancing income distribution and maintaining consumer confidence across various economic cycles.
\\ \textit{Note:} correct because it illustrates how a progressive income tax system serves as an automatic stabilizer that adjusts tax contributions based on income fluctuations, thereby moderating inflation during economic growth and enhancing consumer spending during downturns, which collectively supports overall economic stability.
\end{enumerate}

\vfill
\noindent\textit{End of Answer Key}
\end{document}